\section{Related Literature}\label{sec:literature}

Learning is often introduced in asset pricing models to account for empirical features such as excess volatility or return predictability that are difficult to reconcile with fully informed agents.
The central premise of this literature is that investors have incomplete knowledge of the data generating process governing market outcomes and therefore infer it from observed data.
Our model relates to this large literature on asset pricing with learning. 

Early contributions assume fully rational agents who estimate unknown parameters of a correctly specified model using least squares \citep{Marcet1989}. 
In this setting, learning converges to the rational expectations equilibrium as agents accumulate information over time. 
Even in such low-dimensional environments, learning can generate excess volatility and return predictability \citep{Timmermann1993,Timmermann1996}. 
Convergence to the rational expectations equilibrium is no longer guaranteed when agents are boundedly rational, as finite-memory and finite-sample learning can lead beliefs to fail to converge \citep{Honkapohja2003}. 
The implications of learning become richer when agents face model uncertainty in addition to parameter uncertainty. 
\citet{barberis1998model} study investors who switch between competing forecasting rules, while \citet{hong2007simple} and \citet{Branch2010} analyze settings in which agents estimate univariate perceived laws of motion although the true data generating process is multivariate.

More recently, \citet{adam2011internal} introduce the concept of \textit{internal rationality}, whereby agents choose actions optimally given their subjective beliefs, even when those beliefs are misspecified. 
Building on this insight, \citet{Adam2016} embed internally rational agents into a standard asset-pricing framework in which investors learn the mapping from fundamentals to prices, generating belief--price feedback that helps rationalize key asset-pricing moments. 
We adopt this framework and its central feedback mechanism, but replace their general learning rule with a microfounded LASSO-based learning rule. 
This modification allows us to study how high-dimensional statistical learning and endogenous model selection shape equilibrium price dynamics.

More broadly, our paper also relates to research on how the growth of data affects asset prices. 
\citet{begenau2018big} show that greater data-processing capacity improves investor forecasts and alters equilibrium prices and the cost of capital, while \citet{farboodi2020long} study investors who optimally choose whether to acquire information about fundamentals or about the behavior of other market participants.

On the empirical side, our work relates to the extensive literature on time-varying stock return predictability. 
A wide range of variables, including consumption growth and valuation ratios, forecast returns in some periods but not in others \citep{AngBekaert2007RFS,GoyalWelch2003MS,LettauLudvigson2001JF}. 
A complementary line of research emphasizes sparsity.
Among the many predictors proposed in the literature, only a small and time-varying subset is informative at any given point in time \citep{kozak2020shrinking,freyberger2020dissecting}. 
Consistent with this view, \citet{Mclean2016} document that many return anomalies attenuate following publication, suggesting that investors actively search for, learn about, and trade on such signals.
We add to this empirical literature by providing a theoretical mechanism for the emergence and disappearance  of sparse return predictability.

Finally, our paper contributes to the growing literature on the use of machine learning in finance. 
This literature has largely focused on the empirical gains from machine-learning methods for return forecasting and portfolio choice \citep{Feng2020}. 
\citet{Gu2020} show that machine-learning forecasts outperform traditional linear models by exploiting sparsity and nonlinear interactions among predictors. 
Most closely related to our learning mechanism, \citet{Chinco2019} document that LASSO uncovers short-lived pockets of high-frequency predictability concentrated in a small and shifting set of predictors. 
Our contribution is to shift the focus from forecasting performance to the equilibrium consequences of investors using machine learning techniques in their trading decisions.

% Finally, our work connects to recent efforts to model market outcomes when agents engage in statistical learning.
% \citet{Georges2021} study adaptive model selection and regularization in an agent based market simulation.
% In contrast to this simulation based evidence, our contribution is to provide an analytical characterization of LASSO based learning in a general equilibrium framework and to derive sharp conditions under which sparse and episodic return predictability emerges endogenously.
% --------------------
